\begin{ex}
Tính: $\frac{2}{5} \times \frac{3}{4} = \dots$ (tối giản).
\shortans{3/10}
\loigiai{
$\frac{6}{20} = \frac{3}{10}$.
}
\end{ex}

\begin{ex}
Tính: $\frac{3}{8} : \frac{1}{4} = \dots$
\shortans{3/2}
\loigiai{
$\frac{3}{8} \times 4 = \frac{12}{8} = \frac{3}{2}$.
}
\end{ex}

\begin{ex}
Tính: $\frac{5}{7} \times 14 = \dots$
\shortans{10}
\loigiai{
$14 : 7 \times 5 = 10$.
}
\end{ex}

\begin{ex}
Tính: $15 : \frac{3}{5} = \dots$
\shortans{25}
\loigiai{
$15 \times \frac{5}{3} = 25$.
}
\end{ex}

\begin{ex}
Tính giá trị biểu thức: $\frac{1}{2} + \frac{1}{3} \times \frac{3}{4}$.
\shortans{3/4}
\loigiai{
$\frac{1}{2} + \frac{1}{4} = \frac{3}{4}$.
}
\end{ex}

\begin{ex}
Tính giá trị biểu thức: $(\frac{2}{3} + \frac{1}{3}) : \frac{1}{2}$.
\shortans{2}
\loigiai{
$1 : \frac{1}{2} = 2$.
}
\end{ex}

\begin{ex}
Tìm $\frac{4}{5}$ của 60.
\shortans{48}
\loigiai{
$60 \times \frac{4}{5} = 48$.
}
\end{ex}

\begin{ex}
Một mảnh vườn hình chữ nhật có chiều dài 40 m, chiều rộng bằng $\frac{3}{5}$ chiều dài. Chiều rộng mảnh vườn là bao nhiêu mét?
\shortans{24}
\loigiai{
$40 \times \frac{3}{5} = 24$ m.
}
\end{ex}

\begin{ex}
Chu vi mảnh vườn trong bài toán trên là bao nhiêu mét?
\shortans{128}
\loigiai{
$(40 + 24) \times 2 = 128$ m.
}
\end{ex}

\begin{ex}
Diện tích mảnh vườn trong bài toán trên là bao nhiêu mét vuông?
\shortans{960}
\loigiai{
$40 \times 24 = 960\text{ m}^2$.
}
\end{ex}

\begin{ex}
Tìm $x$: $x \times \frac{2}{5} = 10$.
\shortans{25}
\loigiai{
$x = 10 : \frac{2}{5} = 25$.
}
\end{ex}

\begin{ex}
Tìm $x$: $x : \frac{3}{4} = \frac{8}{9}$.
\shortans{2/3}
\loigiai{
$x = \frac{8}{9} \times \frac{3}{4} = \frac{24}{36} = \frac{2}{3}$.
}
\end{ex}

\begin{ex}
Một cửa hàng có 90 kg đường, đã bán được $\frac{2}{3}$ số đường. Cửa hàng còn lại bao nhiêu ki-lô-gam đường?
\shortans{30}
\loigiai{
Đã bán: $90 \times \frac{2}{3} = 60$ kg. Còn lại: $90 - 60 = 30$ kg.
}
\end{ex}

\begin{ex}
Tính nhanh: $\frac{2}{7} \times \frac{3}{5} + \frac{2}{7} \times \frac{2}{5}$.
\shortans{2/7}
\loigiai{
$\frac{2}{7} \times (\frac{3}{5} + \frac{2}{5}) = \frac{2}{7} \times 1 = \frac{2}{7}$.
}
\end{ex}

\begin{ex}
Tính: $\frac{4}{9} \times \frac{9}{4} = \dots$
\shortans{1}
\loigiai{
1.
}
\end{ex}

\begin{ex}
Tính: $\frac{5}{6} : \frac{5}{6} = \dots$
\shortans{1}
\loigiai{
1.
}
\end{ex}

\begin{ex}
Tính: $\frac{7}{12} \times 0 = \dots$
\shortans{0}
\loigiai{
0.
}
\end{ex}

\begin{ex}
Rút gọn rồi tính: $\frac{6}{9} \times \frac{10}{15}$.
\shortans{4/9}
\loigiai{
$\frac{2}{3} \times \frac{2}{3} = \frac{4}{9}$.
}
\end{ex}

\begin{ex}
Một tấm bìa hình vuông có cạnh dài $\frac{3}{5}$ m. Chu vi tấm bìa là bao nhiêu mét?
\shortans{12/5}
\loigiai{
$\frac{3}{5} \times 4 = \frac{12}{5}$ m.
}
\end{ex}

\begin{ex}
Diện tích tấm bìa hình vuông đó là bao nhiêu mét vuông?
\shortans{9/25}
\loigiai{
$\frac{3}{5} \times \frac{3}{5} = \frac{9}{25}\text{ m}^2$.
}
\end{ex}

\begin{ex}
Tính: $\frac{8}{15} : \frac{4}{5} = \dots$ (tối giản).
\shortans{2/3}
\loigiai{
$\frac{8}{15} \times \frac{5}{4} = \frac{2}{3}$.
}
\end{ex}

\begin{ex}
Tính: $\frac{5}{8} \times \frac{8}{15} = \dots$ (tối giản).
\shortans{1/3}
\loigiai{
$\frac{1}{3}$.
}
\end{ex}

\begin{ex}
Phân số đảo ngược của $\frac{5}{9}$ là phân số nào?
\shortans{9/5}
\loigiai{
$\frac{9}{5}$.
}
\end{ex}

\begin{ex}
Tính: $1 : \frac{2}{7} = \dots$
\shortans{7/2}
\loigiai{
$\frac{7}{2}$.
}
\end{ex}

\begin{ex}
Tìm $\frac{3}{7}$ của 70.
\shortans{30}
\loigiai{
$70 \times \frac{3}{7} = 30$.
}
\end{ex}
